\documentclass[11pt,a4paper]{article}
\usepackage[utf8]{inputenc}
\usepackage[T1]{fontenc}
\usepackage[margin=1in]{geometry}
\usepackage{xcolor}
\usepackage{hyperref}
\usepackage{setspace}
\setstretch{1.15}

\hypersetup{colorlinks=true,linkcolor=blue!70!black,citecolor=blue!70!black,urlcolor=blue!70!black}

\title{\textbf{Point-by-Point Response to Technical Check}\\[0.3em]
\large \textit{npj Artificial Intelligence}}
\author{\textbf{Manuscript ID:} 5a3d24b0-5398-4e89-835d-34efbe8bd9e3 v2.1\\
\textbf{Title:} Financial Time-Series Generation: A Systematic Survey of Taxonomy, Evaluation, Datasets, and Open Challenges\\
\textbf{Authors:} Yingxiao Zhang, Jiaxin Duan, Yue Zou, Ke Feng, Jing Xiao\\
\textbf{Assistant Editor:} Dr.~Monika Samra}
\date{\today}

\begin{document}

\maketitle

\noindent Dear Dr.~Monika Samra and Editorial Team of \textit{npj Artificial Intelligence},

\vspace{0.8em}
\noindent We sincerely thank the editorial team for the careful technical check of our manuscript. We have addressed each of the seven points raised in the Technical Check letter of 7 September 2026. In accordance with your instruction to change nothing beyond the requested items, all revisions are strictly limited to the points below. A marked-up manuscript showing every change in red font is provided under `Related Files' (\texttt{is26\_marked.pdf}), and the clean revised manuscript is provided as \texttt{is26.pdf}. Both compile with zero errors and zero warnings.

\vspace{0.8em}
\noindent Our point-by-point responses follow. Each editor comment is reproduced verbatim in italics, followed by our response and the location of the corresponding change in the revised manuscript.

\vspace{1.5em}
\hrule
\vspace{1.5em}

\subsection*{Point 1: Discussion must not be structured with subheadings}

\vspace{0.5em}
\noindent\textbf{Editor's Comment:}
\textit{``We do not allow a structured Discussion, with subheadings. It can only be divided into paragraphs.''}

\vspace{0.5em}
\noindent\textbf{Response:}
We have removed all subheadings from the Discussion section. In the previous version, the Discussion contained four bolded run-in headings (``Causal and invariant representation learning.'', ``Econometrically bounded generation.'', ``Multi-modal time series generation.'', and ``Online adaptation under regime shifts.''). These bolded lead-ins have been deleted and the text has been rewritten so that each former subheading is now expressed as the opening sentence of a paragraph (e.g., ``A first fundamental challenge concerns causal and invariant representation learning\ldots''). The Discussion (pp.~29--31 of the revised manuscript) now consists of six continuous paragraphs with no subheadings, no bolded lead-in text, and no numbered items, integrating the challenges, limitations, and concluding outlook as required.

\vspace{1.5em}
\hrule
\vspace{1.5em}

\subsection*{Point 2: Results and Methods must use only primary subheadings}

\vspace{0.5em}
\noindent\textbf{Editor's Comment:}
\textit{``Please note that Results and Methods sections should be structured using only primary subheadings. We do not permit the use of secondary subheadings. Any bolded or numbered text placed under primary headings, or any other subsection made under primary subheadings are considered as secondary subheadings. Please edit accordingly.''}

\vspace{0.5em}
\noindent\textbf{Response:}
We have audited the entire Results and Methods sections and removed all secondary-level structure:

\begin{itemize}
  \item \textbf{Results.} The Results section now contains only the seven primary subheadings required by the IRDM structure (Problem Formulation and A Two-Level Taxonomy; Financial Time-Series Extrapolation; Financial Time-Series Imputation; Financial Time-Series Synthesis; Applications; Evaluation Protocols; Data Resources and Benchmarks). No \texttt{\textbackslash subsubsection} commands remain anywhere in the manuscript.
  \item \textbf{Bolded run-in headings removed.} In the Applications subsection, four bolded paragraph headings (``Strategy backtesting and replay.'', ``Stress testing, scenario generation and risk management.'', ``Trading-agent training environments.'', ``Privacy-preserving data sharing and regulatory scenarios.'') have been converted into ordinary sentence openings within their paragraphs (e.g., ``In strategy backtesting and replay, ABIDES\ldots''). No bolded or numbered lead-in text remains under any primary heading.
  \item \textbf{Methods.} The Methods section retains only its four primary subheadings (Literature Retrieval Strategy; Inclusion and Exclusion Criteria; Data Extraction and Coding Protocol; Review Dimensions and Taxonomy Synthesis).
\end{itemize}

\noindent The remaining numbered and bulleted lists in Results and Methods are inline enumerations of items within a paragraph's argument (e.g., the three task-level definitions), not headings: they carry no section numbering, do not appear in the table of contents, and do not interrupt the primary-heading hierarchy. We have verified in the compiled PDF that no bolded standalone heading text appears under any primary subheading.

\vspace{1.5em}
\hrule
\vspace{1.5em}

\subsection*{Point 3: Data Availability Statement must follow the required template}

\vspace{0.5em}
\noindent\textbf{Editor's Comment:}
\textit{``As per our new policies, the Data Availability Statement needs more information. The authors should provide the statement in the following manner: The datasets generated and/or analyzed during the current study are not publicly available due to [INSERT REASON WHY DATA ARE NOT PUBLIC] but are available from the corresponding author on reasonable request. [this is basically for the literature retrieval query logs and systematic extraction coding point mentioned in data availability statement]''}

\vspace{0.5em}
\noindent\textbf{Response:}
We have rewritten the Data Availability statement to follow the required template exactly, inserting the specific reason for the literature retrieval query logs and systematic extraction coding tables. The statement now reads:

\vspace{0.5em}
\noindent\begin{quote}
\itshape ``The datasets generated and/or analyzed during the current study are not publicly available due to the fact that the literature retrieval query logs and the systematic extraction coding tables contain internal reviewer annotations and are subject to licensing restrictions on the underlying subscription databases, but they are available from the corresponding author on reasonable request.''
\end{quote}

\vspace{0.5em}
\noindent The remainder of the statement clarifying that all benchmark datasets, curated databases, and public APIs analyzed in this survey (FinTSB, TFB, TSGBench, CTBench, FinTSBridge, WRDS, CRSP, and FRED) are derived from published literature and publicly accessible repositories has been retained. The revised statement appears in the Declarations section (pp.~32--33 of the revised manuscript). The identical wording has been prepared for the Data Availability field in the submission system (see Point 6).

\vspace{1.5em}
\hrule
\vspace{1.5em}

\subsection*{Point 4: References must follow the standard Nature format}

\vspace{0.5em}
\noindent\textbf{Editor's Comment:}
\textit{``Please ensure the references are in the standard Nature format. References should be numbered sequentially as they appear in the text, methods, tables, and figure legends using a 1, 2, 3 format. They should follow the sequence: author list, title of the paper, name of the journal, volume number, initial-final page numbers or article number (year). Please note that DOIs are required only for online-only publications and correct journal abbreviations should be given. For LaTeX files or manuscripts where bibliography software has been used, please check that the reference list and citations have not been corrupted or display errors. Please either provide a separate BIB/BBL file, or recompile your LaTeX file to include this information.''}

\vspace{0.5em}
\noindent\textbf{Response:}
We have verified and, where necessary, corrected every reference against the Nature format requirements:

\begin{itemize}
  \item \textbf{Sequential numbering.} The bibliography is generated with \texttt{sn-nature.bst} (the Springer Nature reference style) and contains 129 references numbered 1--129 strictly in order of first appearance in the text, Methods, tables, and figure legends. We have programmatically verified that the numbering is sequential with no gaps or duplicates.
  \item \textbf{Element order.} Each entry follows the sequence author list $\rightarrow$ title $\rightarrow$ journal $\rightarrow$ volume $\rightarrow$ pages/article number (year), e.g., ``Cont, R. Empirical properties of trading: a review of the literature. \textit{Quant. Finance} \textbf{1}, 223--236 (2001).''
  \item \textbf{DOIs only for online-only publications.} Print journal articles carry no DOI. Online-only publications (arXiv preprints) are rendered as ``Title. Preprint at \url{https://doi.org/10.48550/arXiv.XXXX.XXXXX} (year).''; 31 such entries appear in the list.
  \item \textbf{Journal abbreviations.} All journal titles use correct standard abbreviations (e.g., \textit{Quant. Finance}, \textit{Manage. Sci.}, \textit{IEEE Trans. Knowl. Data Eng.}, \textit{J. Financ. Econ.}, \textit{Proc. VLDB Endow.}).
  \item \textbf{No corruption.} The manuscript was fully recompiled (pdf\LaTeX{} $\rightarrow$ Bib\TeX{} $\rightarrow$ pdf\LaTeX{} $\times$ 2) with zero errors and zero warnings; all 129 in-text citations resolve correctly to their numbered entries with no ``author?'' or undefined-citation artifacts.
  \item \textbf{BIB/BBL files.} The bibliography database (\texttt{ultimate\_complete.bib}) and the compiled bibliography (\texttt{is26.bbl}) are provided alongside the LaTeX source for production.
\end{itemize}

\vspace{1.5em}
\hrule
\vspace{1.5em}

\subsection*{Point 5: Marked-up manuscript in red font}

\vspace{0.5em}
\noindent\textbf{Editor's Comment:}
\textit{``Marked-up Manuscript: Please include a marked-up manuscript, showing changes to the text in red font. Please do not include tracked changes in comments as we're unable to view them and upload it under `Related Files' section in pdf or docx. Format.''}

\vspace{0.5em}
\noindent\textbf{Response:}
A marked-up manuscript (\texttt{is26\_marked.pdf}, 51 pages) has been prepared and will be uploaded under the `Related Files' section in PDF format. All text changed in this revision round is displayed in red font via a dedicated \texttt{\textbackslash revadd} markup command; no tracked changes or comments are used. The marked-up version contains a summary box on its first page listing the revisions, and every change described in Points 1--3 above (the rewritten Discussion paragraph lead-ins, the converted Applications paragraph openings, and the templated Data Availability statement) appears in red. The clean version (\texttt{is26.pdf}) contains the identical text without color markup.

\vspace{1.5em}
\hrule
\vspace{1.5em}

\subsection*{Point 6: Data Availability in the system must match the manuscript exactly}

\vspace{0.5em}
\noindent\textbf{Editor's Comment:}
\textit{``Please note that The Data Availability information uploaded in the system should match exactly with the statement included in the manuscript.''}

\vspace{0.5em}
\noindent\textbf{Response:}
The Data Availability statement quoted in our response to Point 3 has been copied verbatim into the Data Availability field of the submission system, character for character, so that the system entry matches the manuscript exactly.

\vspace{1.5em}
\hrule
\vspace{1.5em}

\subsection*{Point 7: Author Contributions must match the system exactly, including names}

\vspace{0.5em}
\noindent\textbf{Editor's Comment:}
\textit{``Please note that Author contribution statement provided in manuscript should exactly match with that uploaded in the system which means even the name should be same''}

\vspace{0.5em}
\noindent\textbf{Response:}
The Author Contributions statement in the manuscript (p.~33) reads:

\vspace{0.5em}
\noindent\begin{quote}
\itshape ``Y. Zhang: Conceptualization, Methodology, Writing -- Original Draft. J. Duan: Validation, Formal analysis, Writing -- Review \& Editing. Y. Zou: Validation, Formal analysis, Writing -- Review \& Editing. K. Feng: Supervision, Project administration. J. Xiao: Supervision, Project administration. All authors reviewed, discussed, and approved the final submitted version of the manuscript.''
\end{quote}

\vspace{0.5em}
\noindent This exact text, including the identical author name forms (Y. Zhang, J. Duan, Y. Zou, K. Feng, J. Xiao), has been entered in the Author Contributions field of the submission system so that the two match exactly. Note that the initials Y.~Zhang and Y.~Zou disambiguate the two authors sharing the surname Zhang/Zou as required.

\vspace{1.5em}
\hrule
\vspace{1.5em}

\noindent We hope that the manuscript now satisfies all technical requirements. We have made no changes other than those requested. Should any further adjustments be needed, we would be glad to address them promptly.

\vspace{1.5em}
\noindent Sincerely,\\[0.5em]
Yingxiao Zhang, on behalf of all co-authors\\
\small School of Software and Microelectronics, Peking University

\end{document}
